\documentclass[12pt]{article}
\input{.house-style/preamble.tex}
\usepackage{forest}
\usepackage{xurl}
\usepackage{float}
\useforestlibrary{linguistics}
\setlength{\headheight}{14pt}
\clubpenalty=10000
\widowpenalty=10000
\AtBeginBibliography{\emergencystretch=1em}
\setcounter{biburlnumpenalty}{100}
% Local display normalization; the shared bibliography remains the source.
\DeclareSourcemap{\maps[datatype=bibtex]{\map{
  \step[fieldsource=journal, match=\regexp{^Glossa:\s+a\s+journal\s+of\s+general\s+linguistics$},
        replace={Glossa: A Journal of General Linguistics}]
}}}
\newcommand{\synnode}[2]{\shortstack{\scriptsize #1\\#2}}
\forestset{head edge/.style={edge={line width=1pt}}, nominal tree/.style={for tree={align=center, parent anchor=south, child anchor=north, l sep=5mm, s sep=4mm, inner sep=1.5pt}}}
\hypersetup{pdftitle={Determinatives as nouns in English},pdfkeywords={determinatives, nouns, lexical categories, noun phrases, English}}
\title{Determinatives as nouns in English}
\author{Brett Reynolds \orcidlink{0000-0003-0073-7195}%
\thanks{Contact: \href{mailto:brett.reynolds@humber.ca}{brett.reynolds@humber.ca}}\\
Humber Polytechnic \& University of Toronto}
\date{Draft, September 2026}
\begin{document}
\maketitle

\begin{abstract}
I argue that English determinatives, including articles, demonstratives, and quantifiers, belong within Noun alongside common nouns, proper nouns, and pronouns. Quantificational common nouns provide the closest comparison: they share complement patterns, number transparency, and restricted dependents with parts of the determinative inventory. The strongest adjectival counterweight connects grade, degree modification, and comparative complementation in four quantifiers. Number and referential contrasts supply further nominal connections. Restricted articles belong to the grouping through their integration into the determinative system.

In \mention{take some apples} and \mention{take some}, \mention{some} belongs to the same category. For the independent use, I favour assigning it Head alone, rather than the combined determiner and Head functions in \textit{CGEL}. This gives determinatives the ordinary noun-phrase structure that their inclusion within Noun makes natural. Neither choice requires the other, but their fit supports the package.
\end{abstract}

\noindent\textbf{Keywords:} determinatives, nouns, lexical categories, noun phrases, English

\section{The question and the alternatives}\label{sec:intro}

What is the categorial relationship among words such as \mention{some}, \mention{me}, \mention{apple}, and \mention{Brett}? I argue that all four are nouns: determinatives form a coordinate subcategory alongside common nouns, proper nouns, and pronouns. I write the superordinate category as Noun.

I adopt the general framework of \textit{The Cambridge grammar of the English language} (\textit{CGEL}; \citealt{huddleston2002}). Unlike \textit{CGEL}, though, I include determinatives within Noun. The claims concern synchronic English lexical categories.\footnote{I treat these categories as language-specific: a categorization supported by another language's grammar doesn't determine the English categorization.}

I use \term{determinative} for the category containing articles, demonstratives, and quantifiers such as \mention{the}, \mention{this}, \mention{some}, \mention{every}, and \mention{many}.\footnote{For a fuller inventory, see the online \href{https://www.cambridge.org/highereducation/api/resources/2AC8DF5CA558F8A6FCDFACD0EAF8EA08}{\enquote{List of determinatives in English}} accompanying \textcite{Huddleston2021}.} I reserve \term{determiner} for the syntactic function within the noun phrase (NP) characteristically performed by phrases headed by these words. This distinction separates what kind of word \mention{some} is from what syntactic relationships its phrase participates in.

Compare \mention{take some apples} with \mention{take some}. The word \mention{some} remains determinative in both. Before \mention{apples}, its phrase functions as determiner; without \mention{apples}, the whole expression functions as object. The issue is how these uses fit the lexical taxonomy and the structure of the noun phrase.

Two choices are involved. The first is whether determinatives belong within Noun. The second concerns their Head function in independent NPs. In \textit{CGEL}, the phrase headed by independent \mention{some} jointly fills determiner and Head, a \term{fusion of functions} \citep[410--412]{huddleston2002}. The alternative gives it Head alone.

A standard tree has unique parentage: each constituent below the root has one parent. Fusion allows branches to converge on a shared constituent with two parents. In the \term{ordinary-Head analysis}, the determinative's projection needs no such join. Figure~\ref{fig:some} in §\ref{sec:head-relations} compares the structures.\footnote{The proposal retains fusion elsewhere, including Mod--Head fusion in \mention{the rich}. The preference for unique parentage concerns the determinative's Head configuration, rather than every construction in the grammar.}

The \term{determinative-noun} or \term{D-noun analysis} combines Noun membership with ordinary headedness. The categorization makes ordinary nominal structure a natural treatment, and that regular structure strengthens the case for the grouping. Neither commitment requires the other: a noun can have a fused Head function, and a separate determinative category (D) can permit ordinary headedness. Table~\ref{tab:accounts} shows the four combinations.\footnote{Another logical alternative posits an unrealized noun as Head. I set that option aside here. \textit{CGEL} discusses the limits of ellipsis as a general account \citep[420--421]{huddleston2002}.}

\begin{table}[H]
\centering\small
\caption{Four combinations of categorization and Head analysis. The proposed package is the first row. All keep \mention{apples} as Head in \mention{some apples}.}\label{tab:accounts}
\begin{tabular}{>{\raggedright\arraybackslash}p{3.6cm}>{\raggedright\arraybackslash}p{4.1cm}>{\raggedright\arraybackslash}p{4.4cm}}
\toprule
Account & Place of determinative & Independent \mention{some} \\
\midrule
D-noun, ordinary Head & A subcategory of Noun & Ordinary Head through nominal projection \\
Separate D, ordinary Head & A primary category alongside Noun & Ordinary Head, licensed for both Noun and D \\
Separate D, fused Head & A primary category alongside Noun & Joint determiner and Head functions \\
D-noun, fused Head & A subcategory of Noun & Joint determiner and Head functions \\
\bottomrule
\end{tabular}
\end{table}

Both Head analyses give the whole independent expression NP status. Their difference concerns its internal structure. The four-way comparison leaves room to assess how a categorization and a structural treatment work together, without treating either as entailing the other.

Lexical categories and syntactic functions cut across one another. An NP can function as determiner, as in \mention{\underline{Kim's} book}, and a determinative-headed phrase can function as modifier, as in \mention{the \underline{many} people} \citep{payne2010,pullummiller2022nps}.

The DP hypothesis of \textcite{abney1987} concerns a different Head relation. Under that hypothesis, D heads expressions such as \mention{some apples}. I retain noun-headed NPs and reject the DP analysis for English, following the arguments of \textcite{pullummiller2022nps} and \textcite{bruening2020nominal}.

Here DP means \term{determinative phrase}, as in \textit{CGEL}. This differs from the DP-hypothesis usage just described.

Category membership doesn't remove lexical restrictions, such as the one applying to \mention{every}. \mention{Every apple} is grammatical, but unlike \mention{some} in \mention{I'll take some}, \mention{every} can't occur independently in \ungram{\mention{I'll take every}}. The rules in §\ref{sec:det-uniform} retain that restriction.

NPs already modify adjectives in \mention{a great deal better} and \mention{heaps worse} \citep[549--550]{huddleston2002}. Attested attributive \mention{heaps better} also challenges a general exclusion of NP modifiers from attributive adjective phrases (AdjPs). Degree uses thus extend the comparison with quantificational common nouns; their actual restrictions must be distinguished before NP projection can be counted as an added cost (§\ref{sec:enough}).

\textit{CGEL} already includes pronouns within Noun on the basis of their phrases' functions, despite differences from common and proper nouns in inflection and dependents \citep[327--328]{huddleston2002}. The inclusion of auxiliaries within Verb supplies a further precedent for preserving distinctive properties within a broader lexical category \citep{pullumwilson1977}. These precedents motivate comparing the whole grammatical profile, rather than requiring every member to display each nominal property.

The comparison weighs three considerations: how broadly a property occurs, how it recurs across constructions, and how specifically it connects the groups. These considerations apply equally to nominal and adjectival connections. A pattern concentrated in a few words can be substantial evidence; its contribution depends on its grammatical relationships as well as its reach.

Earlier accounts connect articles with pronouns or give both nominal structure, at different analytical levels (Appendix~\ref{sec:historical}). Hudson's nested noun taxonomy is considered in §\ref{sec:inheritance}; Spinillo's redistribution of determinatives in §\ref{sec:articles}.

Section~\ref{sec:proposal} compares the grammatical profiles, and §\ref{sec:articles} completes the case for the intended inventory. Section~\ref{sec:evidence} then compares internal structures, which the fragment in §\ref{sec:economy} makes explicit and assesses. The further question of coordinate or nested rank follows in §\ref{sec:inheritance}.

\section{Grounds for a broader Noun category}\label{sec:proposal}

The existing Noun category contains common nouns, proper nouns, and pronouns with markedly different meanings, inflection, and dependents. Ordinary count nouns alone are an inadequate comparison group. The question is how determinatives connect with that broader range, and how those connections compare with their affinities to adjectives.

\subsection{Functions and constructional range}\label{sec:independent}

The external syntax of a phrase concerns the functions it fills in larger constructions. Noun phrases can serve as subjects, objects, and complements of prepositions. Determinatives enter these constructions too, with the lexical restrictions illustrated below.

Independent use is widespread within the determinative inventory. \textit{CGEL} discusses such uses for demonstratives and quantifiers including \mention{some}, \mention{all}, \mention{both}, \mention{many}, \mention{few}, \mention{several}, \mention{each}, \mention{either}, \mention{neither}, \mention{much}, and \mention{enough}, while recording lexical restrictions and the separate forms \mention{no}/\mention{none} \citep[371--372, 410--424]{huddleston2002}. The generalization concerns the availability of independent constructions across a lexical category, not unrestricted acceptability in every sentence frame.\footnote{Material before a determiner falls outside the independent-use comparison. \textit{CGEL} treats \mention{all}/\mention{both} in \mention{all/both the books} as predeterminer modifiers, \mention{quite}/\mention{rather} before \mention{a good idea} as peripheral modifiers, and \mention{such}/exclamative \mention{what} before \mention{a disaster} as adjectives \citep[433--437]{huddleston2002}. The fixed \mention{many a} is a complex determinative restricted to Det function \citep[394]{huddleston2002}. The \mention{half} in \mention{half a cake} is a common noun used as a predeterminer modifier \citep[434]{huddleston2002}. These constructions don't add evidence for independent determinative heads.}

In the constructed examples in (\ref{ex:external}), compare the bracketed NPs as subjects, objects, and complements of prepositions. Assume that a woman named Kim and a group of people are already under discussion. The comparison concerns the phrases' functions; §\ref{sec:evidence} considers their internal structure.

\begin{samepage}

\ea\label{ex:external}
\ea \mention{\textup{[}People\textup{]} left.}\qquad \mention{I see \textup{[}people\textup{]}.}\qquad \mention{with \textup{[}people\textup{]}}
\ex \mention{\textup{[}Kim\textup{]} left.}\qquad \mention{I see \textup{[}Kim\textup{]}.}\qquad \mention{with \textup{[}Kim\textup{]}}
\ex \mention{\textup{[}She\textup{]} left.}\qquad \mention{I see \textup{[}her\textup{]}.}\qquad \mention{with \textup{[}her\textup{]}}
\ex \mention{\textup{[}Some\textup{]} left.}\qquad \mention{I see \textup{[}some\textup{]}.}\qquad \mention{with \textup{[}some\textup{]}}
\z\z

\end{samepage}

Independent uses containing fused-head AdjPs prevent a simple inference from these positions to nounhood. \mention{The rich} and \mention{the poor} can fill nominal argument positions. Comparative and superlative adjectives also head expressions without a generic human interpretation: \textit{CGEL}'s \mention{the most important of her criticisms} is an NP containing a partitive \mention{of}-phrase \citep[332--333, 416--423]{huddleston2002}.

Independent \mention{some} can form a one-word NP, whereas an NP with \mention{rich} as fused head requires the definite article \mention{the} on the generic human reading \citep[417--418]{huddleston2002}. That is a local contrast: argument NPs headed by singular count common nouns also need determination. The broader comparison has to include dependents, inflection, and interpretation as well as external position.

Existentials add a construction to this comparison. With possible solutions under discussion, the constructed \mention{There are several} places the independent quantifier in displaced-subject position; dummy \mention{there} is the subject \citep[1391--1393]{huddleston2002}. Section~\ref{sec:quant-controls} compares this use with quantity words near the adjective--determinative boundary, distinguishing expressions with partitive or relative dependents from those interpreted through context alone.

Determiner function also cuts across the existing noun subcategories. In \mention{my preferences}, \mention{Kim's preferences}, and \mention{people's preferences}, the determiner is an NP ultimately headed by a pronoun, a proper noun, and a common noun respectively \citep[354--355, 470--471]{huddleston2002}.\footnote{\textit{CGEL} assigns these genitives the combined function Subject--Det \citep[472--473]{huddleston2002}. I treat them as Det here, without the additional subject function.}

The compound determinative \mention{someone}, following the categorization of \textcite[§1.3]{Payne2007}, participates in \mention{someone's preferences}. Plain-case NPs also fill Det: \mention{what size} in \mention{what size shoes}, \mention{that size} in \mention{that size shoes}, and \mention{Sunday} in \mention{Sunday morning} \citep[356]{huddleston2002}.

Determiner remains the characteristic function of determinative phrases. This specialization could support a separate primary category. Its tasks are themselves closely connected with nominal reference: Det marks definiteness and often contributes quantification \citep[354--359]{huddleston2002}. The issue is whether that specialization warrants a primary boundary or a distinction within Noun.

The existing NP determiners provide a positive comparison. In \mention{Kim's preferences}, the embedded NP supplies an identifying anchor through its own referent; \mention{Sunday} and \mention{what size} specify a day or a dimension. Determinatives' deictic and quantitative specifications fit this nominal pattern. I take that fit to support treating their specialization in Det as a distinction within Noun.

Preposition phrases (PPs) remain a restricted alternative, as in \mention{up to twenty minutes} and \mention{between fifty and sixty tanks} \citep[356]{huddleston2002}. Number, countability, and other selectional conditions distinguish the determining expressions.

Modifier and adjunct functions also cut across the groups. Compare the modifiers in \mention{dog houses}, \mention{Canada Day}, \mention{the manager herself}, and \mention{the few people}. Emphatic \mention{herself} also functions as a clause adjunct in \mention{The manager detected the error herself} \citep[1496--1497]{huddleston2002}. Temporal NPs such as \mention{that day} and \mention{Sunday} supply adjuncts, as does degree \mention{enough} in \mention{I hadn't prepared enough} (§\ref{sec:enough}). These are shared functions with construction-specific distributions.

Relative-clause postmodification supplies a further constructional comparison. Common nouns freely take integrated relatives, as in \mention{people who came}; personal pronouns permit a restricted range, including \mention{we who have read the report} \citep[430]{huddleston2002}. Determinative examples include \mention{few who come ever leave}, \mention{those who came}, \mention{that which remains}, and \mention{something that you need to know}.

Compounds also permit \mention{anyone who asks} and \mention{everything that matters}. With books under discussion, compare \mention{some that I saw} and \mention{two that I have seen}. \textit{CGEL} describes relative postmodification with demonstratives and compounds \citep[414, 422--423]{huddleston2002}.

Compound determinatives also take post-head adjectives. CGELBank, a treebank annotated in the \textit{CGEL} framework \citep{reynolds2023unified}, attests \mention{I need something reliable and good looking}.\footnote{Sentence \href{https://github.com/nert-nlp/cgel/blob/d0a2c2d8c522b23d5aa879fc9cb94948f4285b2c/datasets/ewt-test_iaa50.cgel\#L179-L181}{\nolinkurl{answers-20111024111513AAAQhAO_ans-0003}} in \nolinkurl{ewt-test_iaa50.cgel}.} The position and interpretation of these modifiers have specialized conditions (§\ref{sec:compounds}).

\subsection{The connection from quantificational common nouns}\label{sec:quant-nouns}

Quantificational common nouns make the connection more specific than shared quantity meanings. In \mention{a lot of the delegates} and \mention{many of the delegates}, both heads quantify over a partitive \term{domain}: the whole from which a quantity is drawn, here the delegates. But \mention{lot} also permits \mention{a lot of delegates}, whereas \ungram{\mention{many of delegates}} is excluded \citep[349]{huddleston2002}. The comparison links the groups through a shared construction while preserving a specific difference in complementation.

Agreement provides another connection. Compare the constructed \mention{A lot of the people were waiting} and \mention{Some of the people were waiting}, with plural agreement, against \mention{A lot of the work was finished} and \mention{Some of the work was finished}, with singular agreement. The whole subject NP's number depends on the complement of \mention{of}. \textit{CGEL} calls quantificational \mention{lot} \term{number-transparent} \citep[349--350, 411--412]{huddleston2002}.

This pattern connects particular constructions, not every member of either category. Singular \mention{each of the people} doesn't take its number from \mention{people}; nor does ordinary \mention{a photograph of the people}. Number transparency strengthens the comparison between quantificational \mention{lot} and number-neutral \mention{some} without turning complement number into a general rule for Noun.

Restricted dependents also occur on the common-noun side. \textit{CGEL} categorizes quantificational \mention{plenty} as a common noun whose use resists determination and modification; \mention{lot} requires \mention{a} and permits only limited modification \citep[349--350]{huddleston2002}. Established common nouns thus approach the determinative profile as determinatives approach theirs. Ordinary count nouns alone provide an inadequate comparison group.

The connection also extends beyond argument NPs. Quantificational nouns occur in degree modifiers such as \mention{a great deal smaller} and \mention{plenty big enough} \citep[549--550]{huddleston2002}. Section~\ref{sec:enough} compares these with determinative degree modifiers, including their attributive distribution. Meaning, complementation, agreement, and restricted dependents give the proposed grouping a specific basis within established Noun.

The comparison is specific to these constructions. Common nouns also permit a wider range of PP and clausal complements, whereas pronouns and primary naming uses of proper nouns have much more restricted dependents \citep[429--430, 439--443, 517--521]{huddleston2002}. Partitives such as \mention{a lot/some of the wine} connect established common nouns with determinatives without making their complement systems identical.

\subsection{Quantificational adjectives and nominal independence}\label{sec:quant-controls}

\mention{Numerous people} and \mention{many people} both quantify, but their syntax differs. Degree modification exposes a contrast: \mention{so many mistakes} is possible, while \ungram{\mention{so numerous mistakes}} is excluded \citep[539--540]{huddleston2002}. Quantity meaning and prenominal position provide a semantic control for the nominal comparison. The question is which further constructions these words enter and under what conditions.

The boundary already cuts through this semantic range in \textit{CGEL}. It includes \mention{several} among determinatives and treats \mention{certain} and \mention{various} as marginal members, partly on the evidence of partitives \citep[392--393]{huddleston2002}. Their relevant readings differ from adjectival \mention{certain} meaning ‘sure’ and the distributive reading of \mention{their several ways}. The comparison below concerns their quantificational uses, alongside \mention{numerous}, \mention{multiple}, and \mention{countless}.

Predicative use supplies another baseline. \textit{CGEL} records \mention{Its advantages are several} and \mention{Their enemies were many}, describing the latter pattern as uncommon and formal \citep[392, 395--396]{huddleston2002}. Compare adjectival \mention{Their enemies were numerous}. Since these uses span the boundary, their value lies in combination with other properties. \mention{Every} and the articles remain excluded from this predicative pattern.

Independent constructions differ in how their interpretation is supplied. In \mention{certain of them}, a PP gives the partitive domain; in \mention{many who came}, a relative adds descriptive content. Anaphoric \mention{many} can stand alone. This suggests a possible cline of decreasing dependence on overt nominal support, starting from an ordinary noun such as \mention{people}. Partitives and relatives are alternative supports whose relative permissiveness remains to be established.

\textit{CGEL}'s restrictions give this hypothesis a starting point. The \mention{of}-phrase with independent \mention{certain} and \mention{various} is hardly omissible, whereas independent \mention{many} and \mention{several} allow both explicit partitives and anaphoric uses without the PP \citep[392, 411--413]{huddleston2002}. The partitive with \mention{certain} is relatively formal; that with \mention{various} is accepted by some speakers, primarily in American English \citep[393]{huddleston2002}.

Table~\ref{tab:quant-controls} separates NP dependents from the surrounding construction. The argument frames are \mention{X of them}, \mention{X who arrived left}, and anaphoric \mention{X left}. The existential frames are \mention{There are X of them}, \mention{There are X who can help}, and anaphoric \mention{Are there X?} Interpret the constructed comparisons with people under discussion. The common-noun controls retain \mention{a} in \mention{a lot}.

\begin{table}[H]
\centering\small
\caption{Independent quantificational uses. G: described in \textit{CGEL}; R: restricted there; I: constructed positive illustration in the stated frame. Numbered references identify attestations below. A question mark marks an unresolved cell, not an ungrammaticality judgment.}\label{tab:quant-controls}
\setlength{\tabcolsep}{3pt}
\begin{tabular}{>{\raggedright\arraybackslash}p{1.65cm}*{6}{>{\centering\arraybackslash}p{1.68cm}}}
\toprule
 & \multicolumn{3}{c}{Argument use} & \multicolumn{3}{c}{Existential displaced subject} \\
\cmidrule(lr){2-4}\cmidrule(lr){5-7}
Expression & Partitive PP & Relative clause & No dependent & Partitive PP & Relative clause & No dependent \\
\midrule
\mention{many} & G & G & G & I & I & I \\
\mention{several} & G & I & G & I & I & I \\
\mention{certain} & G & ? & R & ? & (\ref{ex:certain-exist}) & ? \\
\mention{various} & G & ? & R & ? & ? & ? \\
\mention{numerous} & (\ref{ex:numerous-part}) & ? & (\ref{ex:numerous-subj}) & ? & (\ref{ex:numerous-exist}) & ? \\
\mention{multiple} & (\ref{ex:multiple-part}) & ? & (\ref{ex:multiple-subj}) & (\ref{ex:multiple-exist-part}) & (\ref{ex:multiple-exist-rel}) & (\ref{ex:multiple-bare}) \\
\mention{countless} & (\ref{ex:countless-part}) & ? & ? & (\ref{ex:countless-exist-part}) & (\ref{ex:countless-exist-rel}) & ? \\
\mention{plenty} & G & I & I & I & I & I \\
\mention{a lot} & G & I & I & I & I & I \\
\bottomrule
\end{tabular}
\end{table}

The G and R cells retain the qualifications just stated; the common-noun partitives follow §\ref{sec:quant-nouns}. The relative use of \mention{many} follows \textit{CGEL}'s discussion of the human interpretation \citep[414]{huddleston2002}. The attestations below combine targeted searches in NOW and COCA with separately located examples. Expanded contexts exclude overt following nouns, predicative inversions, and mathematical \mention{multiple}. These searches establish selected occurrences; they supply no usage rates.\footnote{Searches conducted on 12 September 2026 in \href{https://www.english-corpora.org/now/}{NOW} and \href{https://www.english-corpora.org/coca/}{COCA}, using surface strings without part-of-speech filters. Queries, retained excerpts, and verification limits are recorded in \nolinkurl{notes/2026-09-12-quantifier-controls-corpus.json}; separate source checks are in \nolinkurl{notes/2026-09-12-quantifier-controls-sources.json}.}

The examples in (\ref{ex:quant-arguments}) show partitive and independent subject uses of further quantity words. The \mention{countless} partitive comes from academic prose; the \mention{multiple} partitive from a blog discussion. The final two examples have independent quantifiers as subjects, with people recoverable from context. Their occurrence requires a description even where speakers differ over their acceptability.

\ea\label{ex:quant-arguments}
\ea\label{ex:countless-part} \mention{countless of their particular statements and deeds}\footnote{Thomas Hughson, \href{https://theologicalstudies.net/wp-content/uploads/2022/08/69.1.1.pdf\#page=7}{\enquote{Interpreting Vatican II: ‘A New Pentecost’}}, \textit{Theological Studies} 69 (2008), p.~9; COCA academic record, checked against the journal PDF. The NP is object of \mention{irradiate}.}
\ex\label{ex:numerous-part} \mention{numerous of its inhabitants}\footnote{Norman Harrington, \textit{New York Times}, 7 April 1968, as quoted in \textit{Merriam-Webster's Dictionary of English Usage} and reproduced in Reynolds and Pullum's \href{https://www.lel.ed.ac.uk/~gpullum/ClosedClasses.pdf}{\enquote{New members of closed classes in English}}, §2.1.2. Verified through that draft; the newspaper original wasn't checked.}
\ex\label{ex:multiple-part} \mention{multiple of the fast casual eateries}\footnote{COCA, 2012 blog discussion under \href{http://www.tonetoatl.com/2012/08/burgerfi-bringing-its-burgers-to.html}{\enquote{BurgerFi Bringing Its Burgers to\ldots}}. The surrounding phrase is \mention{in front of multiple of the fast casual eateries}. Expanded corpus context and source metadata checked; the original page wasn't available.}
\begin{samepage}
\ex\label{ex:numerous-subj} \mention{numerous were injured}\footnote{\href{https://pagesix.com/2022/07/04/travis-scott-performs-in-miami-at-e11even/}{Page Six}, 4 July 2022; NOW expanded context. The clause follows a coordinated clause about people who died. Here \mention{numerous} is the canonical subject.}
\ex\label{ex:multiple-subj} \mention{multiple were injured}\footnote{Ivana Saric and Herb Scribner, \href{https://www.axios.com/2022/06/27/amtrak-train-derail-missouri}{Axios}, 27 June 2022, opening sentence; original page checked.}
\end{samepage}
\z\z

Existentials show a further range in (\ref{ex:quant-existentials}). In (\ref{ex:multiple-bare}), \mention{multiple} ends the question, with wedding dresses under discussion. The next two examples contain partitives. The remaining examples have relative postmodifiers and include published prose, a hearing transcript, and reported speech. Source locations identify the records; they don't establish each speaker's or writer's variety.

\ea\label{ex:quant-existentials}
\ea\label{ex:multiple-bare} \mention{are there multiple?}\footnote{Kira Bindrim, in a \href{https://qz.com/india/2130303/the-big-fat-indian-wedding-is-only-getting-bigger}{Quartz podcast transcript}, 2022; NOW and original transcript. The preceding question introduces wedding dresses. NOW dates the record 22 February; the current page reports an update on 20 July.}
\ex\label{ex:multiple-exist-part} \mention{there are multiple of them}\footnote{Ivan Mehta, \href{https://techcrunch.com/2022/09/16/ios-16-users-are-getting-creative-with-the-new-background-removal-feature-for-photos/}{TechCrunch}, 16 September 2022; NOW and original article. The antecedent is objects in a photo.}
\ex\label{ex:countless-exist-part} \mention{There are countless of them out there}\footnote{\href{https://seekingalpha.com/article/4349886-facebooks-future-is-brighter-ever}{\enquote{Facebook's Future Is Brighter Than Ever}}, Seeking Alpha, 25 May 2020; NOW expanded context and source metadata. The antecedent is payment apps.}
\ex\label{ex:multiple-exist-rel} \mention{There are multiple that are cited in my response.}\footnote{Zachary Loyed, public-records hearing, \textit{State of Florida v.~Thomas Lee Gudinas}, 29 May 2025, transcript p.~7, lines 14--15; \href{https://www.supremecourt.gov/DocketPDF/24/24-7457/363457/20250618100158557_Appendix\%20to\%20the\%20Petition\%20for\%20Writ\%20of\%20Certiorari-06-18-25.pdf}{Appendix C}, PDF p.~55. Original transcript checked; \mention{multiple} refers to cases.}
\ex\label{ex:countless-exist-rel} \mention{There are countless who have greatly impacted their communities}\footnote{\href{https://metro.co.uk/2023/10/01/the-sisterhood-list-saluting-the-uk-black-women-whove-opened-doors-for-others-19558944/}{\enquote{The Sisterhood List}}, Metro, 1 October 2023; NOW expanded context and source metadata.}
\ex\label{ex:numerous-exist} \mention{there are numerous who are willing to donate}\footnote{Nana Patekar, as quoted in English by IANS in \href{https://indianexpress.com/article/entertainment/bollywood/nana-patekars-foundation-collects-rs-80-lakh-for-drought-hit-farmers/}{The Indian Express}, 19 September 2015; NOW and original article. The original language of the speech wasn't established.}
\ex\label{ex:certain-exist} \mention{there are certain who would not like} \ldots\footnote{Donald Trump, as quoted in \href{https://www.foxnews.com/politics/trump-throws-support-behind-coronavirus-stimulus-bill-in-senate-vows-to-sign-vital-legislation-immediately}{Fox News}, 25 March 2020; NOW record dated 26 March and original article. The quotation was checked against the article, without an audio check.}
\z\z

These examples motivate comparing each expression with and without its dependent. They don't establish that adding a relative improves acceptability. A cline would require matched judgments showing that use with less overt support predicts use in the more supported frames, for the same reading and speaker population. Crossings would count against a single ordering. Canonical argument and existential frames remain separate comparisons; no historical sequence follows from their present distribution.

The ordering would also be local to the quantificational uses. Independent genitives exclude direct partitives, while partitive \mention{those} requires a following modifier \citep[413]{huddleston2002}. Special human interpretations, such as \mention{Few would disagree}, add a lexical distinction beyond anaphora \citep[414]{huddleston2002}. Nor does the obligatory article in quantificational \mention{a lot} make its common-noun head less nominal. The cline concerns conditions on independent use, rather than a general scale of nounhood.

The existential comparison narrows the contrast with generic \mention{the rich}, but discourse conditions matter. \textit{CGEL} allows definite displaced NPs in several discourse configurations \citep[1397--1401]{huddleston2002}. Adjective fusion also permits indefinite ordinal \mention{a second} \citep[416]{huddleston2002}, giving the constructed comparison \mention{There was a second}. Article-free use distinguishes some quantity words within this broader range; existential occurrence alone leaves fusion available.

Agreement likewise belongs initially to the phrase. In \textit{CGEL}'s account, \mention{there} inherits person--number from the displaced NP, with informal invariant \mention{there's} as a complication \citep[242]{huddleston2002}. Plural agreement also occurs with adjective-containing \mention{the rich}, whose adjective lacks number inflection \citep[418]{huddleston2002}. The number transparency of \mention{some} and \mention{a lot} in §\ref{sec:quant-nouns} supplies the more specific comparison, with explicit count and non-count domains.

The controls sharpen what the broader Noun proposal has to explain. Quantity words share some independent constructions while differing in degree syntax, determination, and the conditions on their dependents. Categorizing every new independent use as determinative solely because it passes this test would make the comparison circular. I assess Noun membership through the recurring combination of properties, with the constructional restrictions retained.

\subsection{The connected adjectival profile}\label{sec:adjectival-profile}

The strongest adjectival connection also concerns a restricted group: \mention{few}, \mention{many}, \mention{much}, and \mention{little}. They distinguish plain, comparative, and superlative forms: \mention{few}/\allowbreak\mention{fewer}/\allowbreak\mention{fewest}, \mention{many}/\allowbreak\mention{more}/\allowbreak\mention{most}, and corresponding paradigms for \mention{little} and \mention{much} \citep[391--395]{huddleston2002}. Grade combines with degree-modifier selection and comparative complementation, as in \mention{more than ten}. The three properties form a connected profile.

Degree \mention{very} selects gradable heads: \mention{very tall} and \mention{very few}, but \ungram{\mention{very you}}, \ungram{\mention{very the book}}, and \ungram{\mention{very my life}}. It also excludes non-gradable determinatives: \ungram{\mention{very every}}, \ungram{\mention{very some}}, and \ungram{\mention{very this}}. The associated series includes \mention{so few}, \mention{too few}, \mention{how few}, and \mention{as few as}. These degree patterns select \mention{few}, \mention{many}, \mention{much}, and \mention{little}, the four determinatives that inflect for grade \citep[393--395, 431--432]{huddleston2002}.\footnote{The restriction concerns this degree series. \textit{CGEL} also records modifiers of completeness, including \mention{marginally enough} and \mention{absolutely all}, with different selectional ranges \citep[432]{huddleston2002}. Those combinations don't extend the \mention{very}/\mention{so}/\mention{too}/\mention{how} series beyond the four degree determinatives.}

With the complex determinative \mention{a few}, modifier position distinguishes two kinds of attachment. \textit{CGEL} explicitly contrasts internal \mention{very} in \mention{a very few mistakes} with peripheral \mention{quite} in \mention{quite a few mistakes} \citep[p.~392, (61)]{huddleston2002}. The reversed orders \ungram{\mention{very a few}} and \ungram{\mention{a quite few}} are excluded in these constructions. Here \mention{a} belongs to the complex determinative; \mention{very} occurs within it and \mention{quite} outside it.

Approximative and focusing modifiers have a broader distribution. Compare \mention{almost every}, \mention{nearly all}, \mention{hardly any}, \mention{practically no}, \mention{almost everyone}, and \mention{almost ten}. They have counterparts in established NP modification: \mention{almost the whole class}, \mention{hardly a day}, and \mention{only the best}. In \mention{almost every student}, \mention{almost} approximates the quantity; it doesn't intensify a gradable property of \mention{every}.

The broader approximative pattern supplies NP parallels distinct from the four quantifiers' adjectival profile. Section~\ref{sec:payne} develops an attachment analysis of that contrast, available under either ordinary-Head taxonomy.

Overlap with adjectives likewise needs phrase levels kept distinct. Both \mention{the people} and \mention{the rich} are NPs in \textit{CGEL}, though only the former has a noun as lexical head. In \mention{a soccer/round ball}, a noun-headed expression and an AdjP share modifier function; in \mention{Jones became president/ill}, NP and AdjP share predicative complement function. External functions alone therefore don't establish nounhood.

\mention{She is a beauty} ascribes a property, whereas \mention{That is Kim} identifies a person. \textit{CGEL} excludes specifying \mention{be} clauses from its noun--adjective diagnostic because phrases from other categories occur in them too \citep[536]{huddleston2002}.

Nominal grammar also distinguishes predicative and argument uses. A bare-role NP such as \mention{president} is licensed in \mention{I'd like to be president}, but requires determination in the corresponding object use \mention{I'd like to meet the president} \citep[328]{huddleston2002}. Determinatives' asymmetries fit this broader variation within Noun.

Bare colour expressions and evaluative comparative subjects complicate a simple distributional boundary. \textit{CGEL}'s \mention{Henrietta likes red shirts, and I like blue} permits reduction licensed by coordination. The constructed \mention{Blue is good} raises a separate ambiguity between noun and adjective, while the attested \mention{Bluer is better} establishes a bare comparative subject without settling its phrase category.\footnote{The \mention{blue}, \mention{old}, and \mention{small} examples on \textit{CGEL} p.~417 all occur in coordinated contrasts. \mention{Bluer is better} accompanies a results colour scale in Stanford CS329X, \href{https://web.stanford.edu/class/cs329x/slides/Lecture12_B_Codeswitching\%20LLMs.pdf}{\enquote{Codeswitching LLMs}}, slide~8 (accessed 10 September 2026). I interpret it as evaluating a degree of blueness. Subject function alone doesn't establish NP status \citep[236]{huddleston2002}.} These cases leave the wider comparison of argument distribution and head properties necessary.

\subsection{Number, genitive marking, and reference}\label{sec:inflection}

Inflection supplies several connections with established nouns, though no one contrast runs through the whole category. Common nouns typically distinguish singular and plural; proper nouns permit number inflection in restricted uses such as \mention{the Smiths}. Demonstratives likewise distinguish \mention{this}/\mention{these} and \mention{that}/\mention{those} \citep[373, 521]{huddleston2002}. These are contrasts between forms of a lexeme, unlike the fixed number restrictions of \mention{each} and \mention{several}.

Number needn't be inflectional wherever it matters. \textit{CGEL} treats singular and plural \mention{you} as distinct lexemes, differentiated overtly only by the reflexive forms \mention{yourself} and \mention{yourselves} \citep[486]{huddleston2002}. Number-neutral \mention{some} has no corresponding singular--plural alternation. The demonstratives supply positive morphological evidence; the absence of that alternation elsewhere doesn't distinguish determinatives from pronouns as whole categories.

The demonstrative number contrasts remain visible in independent uses. Generic human \mention{the rich}, by contrast, is a plural NP whose adjective lacks number inflection and still takes adverb modifiers: \mention{the very rich} \citep[418]{huddleston2002}. The difference concerns overt paradigms, rather than the mere possession of a number value by the whole NP.

Case gives a further partial connection. Genitives occur across the proposed subcategories: \mention{dog's}, \mention{Kim's}, \mention{my}, and \mention{someone's}. Personal pronouns have fuller case paradigms. Most determinatives lack case inflection, and compounds such as \mention{someone} and \mention{something} may owe their genitives to the nominal component \citep[423--424, 479--480]{huddleston2002}. This supports the compounds' nominal affinity more directly than the category's membership as a whole.

Meaning and reference vary within Noun. \mention{Apple} describes a category; \mention{Kim}, in its primary naming use, identifies through a name; \mention{she} depends on context. Determinatives range from deictic \mention{this} to quantitative \mention{many} and universal \mention{every} \citep[358--360, 370--405, 425--428, 515--521]{huddleston2002}. Deixis connects demonstratives with pronouns, while quantity connects determinatives with common nouns such as \mention{number} and \mention{majority}.

These meanings aren't exclusive to nouns. Adjectives such as \mention{singular} and \mention{plural} also concern number; \mention{proximate} and \mention{distal} describe spatial relations relevant to demonstrative contrasts, though spatial meaning needn't itself be deictic. Nor does Noun imply close semantic similarity throughout: \mention{every} quantifies over a nominal restriction, while \mention{Kim} identifies an individual. Semantic affinities contribute to the comparison without settling its boundaries.

Pro-form gender \citep{reynolds2025proformgender} connects these ways of referring. Descriptions such as \mention{the woman} and names such as \mention{Kim} identify referents whose construal bears on pro-form choice. Pronouns and determinatives express contrasts between personal \mention{she}/\mention{somebody} and non-personal \mention{it}/\mention{something}. Relative \mention{who}/\mention{which} supplies another contrast.\footnote{Whether independent \mention{what} and relative \mention{which} belong to pronoun or determinative leaves them within the proposed Noun category. I leave that internal boundary open here; \textit{CGEL}'s inventory suffices for the present comparison \citep[397--399]{huddleston2002}.}

Interrogative and relative properties are distinct from functions such as object. An interrogative object may be headed by a pronoun (\mention{who}) or a common noun (\mention{which book}); the latter obtains its interrogative property from a dependent. Neither the object function nor the interrogative construction determines the lexical head's category.

Inflection and reference provide partial connections, with different reach. Demonstrative number directly compares forms of determinative lexemes with nominal number paradigms; compound genitives may depend on their nominal component. The \mention{no}/\allowbreak\mention{none} alternation concerns form selection, discussed with restricted membership in §\ref{sec:form-selection}. None of these properties is a necessary condition for every member of Noun.\footnote{Derivations such as \mention{nothingness} and \mention{oneness} are weaker category diagnostics: \mention{-ness} also attaches to adjectives and other bases \citep[Ch.~19, §5.7.2(i)]{huddleston2002}. Numeral morphology likewise needs its inputs distinguished. \textcite[§§5.4--5.5]{reynolds2026numerals} derives fractional nouns from cardinal nouns and treats \mention{two thousand and twenty-seventh} as a coordination whose final coordinate is adjectival. Neither process independently establishes determinative nounhood.}

\subsection{Weighing the profiles}\label{sec:membership}

Table~\ref{tab:existing} summarizes the comparison. Functions concern phrases, while inflection concerns word forms. AdjP and AdvP denote adjective and adverb phrases. The modification rows identify proposed determinative attachments, explained in §\ref{sec:payne}. The rows aren't equally weighted tests, and several properties belong to one connected pattern.

\begin{table}[H]
\centering\small
\caption{Profiles compared. Determinative modifier attachments are analytical commitments; interrogative and relative forms follow \textit{CGEL}'s inventory.}\label{tab:existing}
\setlength{\tabcolsep}{4pt}
\begin{tabular}{>{\raggedright\arraybackslash}p{2.25cm}*{4}{>{\raggedright\arraybackslash}p{2.42cm}}}
\toprule
Dimension & Common noun & Proper noun & Pronoun & Determinative \\
\midrule
Argument / Det functions & NP arguments; genitive Det & NP arguments; genitive Det & NP arguments; genitive Det & NP arguments; plain and genitive Det \\
Accepts determination & Broad contrasts; singular count arguments normally require it & Restricted in primary naming uses & Normally excluded & Lexically restricted: \mention{the few}, \mention{these three} \\
Internal modification & Productive AdjP and nominal premodifiers; relative postmodifiers & Restricted AdjP premodifiers and embellishments & Restricted AdjP premodifiers and relative postmodifiers & Degree AdvP on four quantifiers; NP modifiers; \mention{the lucky few}; relatives \\
Complemen\-tation & Selected PPs and clauses & Not characteristic of primary naming uses & Normally absent & Partitive \mention{of}-PPs; comparative \mention{than}-phrases \\
Peripheral AdvP & \mention{only the book} & \mention{even Kim} & \mention{only you} & \mention{almost every}; \mention{hardly any}; \mention{almost ten} \\
Modifier / adjunct functions & \mention{dog houses}; \mention{that day} & \mention{Canada Day}; \mention{Sunday} & Emphatic reflexives in both functions & \mention{the few people}; degree adjuncts \\
Meaning & Descriptive properties and relations & Naming in primary uses & Person, deixis, anaphora, and interrogation & Quantification, definiteness, deixis, and interrogation \\
Pro-form gender & Descriptions shape referent construal & Names identify referents construed by gender & Gender-sensitive forms & Gender-sensitive forms and constructions \\
Inflection & Number and genitive & Genitive; restricted number & Case, including genitive; reflexive forms & Number; grade; genitive; \mention{no}/\mention{none} \\
Interrogative / relative constructions & Through another constituent, e.g.\ \mention{which book} & Through another constituent & \mention{who}, \mention{whose} & \mention{which}, \mention{what}, and \mention{-ever} forms \\
Inventory & Open & Open to new names & Closed & Closed \\
\bottomrule
\end{tabular}
\end{table}

Both the nominal and adjectival connections are unevenly distributed. Quantificational common nouns supply a connected comparison involving quantity, complements, agreement, and restricted dependents. The four gradable quantifiers supply a connected adjectival profile. Concentration in a subgroup doesn't make either comparison negligible. Their significance depends on how each pattern connects with the rest of the inventory.

The quantity-word controls in Table~\ref{tab:quant-controls} show where the independent constructions overlap. \mention{Certain}, \mention{numerous}, \mention{multiple}, and \mention{countless} have different documented ranges. Independent use supplies one part of their profile. The comparison also retains differences in degree syntax and the conditions on partitives, relatives, and anaphora; the proposed ordering of these constructions remains to be tested.

I give greater weight to the recurring nominal profile. Argument and determiner functions, partitives, relative postmodification, and number and referential contrasts connect determinatives with established parts of Noun. Within this range, the gradable quantifiers retain their determining and independent uses alongside the degree pattern. Keeping them within determinative expresses those connections while locating the grade and degree restrictions in a narrower group.

The broader Noun grouping still needs an account of restricted members. Their inclusion depends on how they participate in the determinative system.

\section{Restricted members and the scope of the grouping}\label{sec:articles}

The category-level argument draws on constructions that many determinatives enter. Articles lack ordinary independent uses, so their inclusion needs a further argument: they belong to the determinative system whose broader profile supports Noun membership. The \mention{no}/\mention{none} paradigm raises a related question about restrictions on forms.

\subsection{Articles within the determinative system}\label{sec:article-membership}

Why categorize the articles \mention{the} and \mention{a} as nouns if they can't stand independently? \mention{Every} is restricted too, while \mention{no} has the distinct independent form \mention{none} \citep[371--372, 410--411]{huddleston2002}.

\textit{CGEL}'s treatment of \mention{my} supplies a precedent for restricted membership within Noun \citep[470--471]{huddleston2002}. Dependent \mention{my} and independent \mention{mine} belong to one pronoun paradigm. \mention{My} remains a noun even though, in the ordinary NP construction, it requires a following nominal. For articles, the corresponding positive evidence is integration into determinative; their nounhood depends on the category-level argument.

The articles participate in the determinative system of definiteness, quantity, and count restrictions. In \mention{the/a/this/every book}, they occupy the same Det position and contribute to the interpretation of the NP. \mention{The} permits singular, plural, and non-count targets; \mention{a} selects a singular count target and contributes individuation. Their connections extend beyond shared position, although neither has the full determinative profile \citep[368--373]{huddleston2002}.

\mention{The} also participates in degree modification. In \mention{the bigger the better}, it modifies comparative AdjPs, as other determinatives do in \mention{much bigger} or \mention{no better} \citep[1131--1132, 1135--1136]{huddleston2002}. This supplies a connection outside ordinary Det function while preserving its exclusion from independent argument use.

\mention{A} also enters the complex determinatives \mention{a few}, \mention{a little}, and \mention{many a}. In the last, \mention{many} contributes a large number and \mention{a} an individuating, distributive effect. \textit{CGEL} treats \mention{many a} as syntactically fixed: it doesn't establish that ordinary \mention{a} freely accepts modifiers \citep[392--394]{huddleston2002}. It does connect the article with quantificational constructions beyond \mention{a book}.

\textcite[153--158]{spinillo2004reconceptualising} instead retains \mention{the}, \mention{a}, and \mention{every} as an expanded article category and redistributes other determinatives among adjectives and pronouns \citep[194--195]{spinillo2004reconceptualising}. Her grounds include dependence on a following nominal, lack of predicative and partitive uses, and limited descriptive content. She also recognizes differences within the article trio. The issue is whether this cluster warrants separating the trio from the wider determinative system.

\mention{Every} connects with independent \mention{each} through universal quantification and singular count selection. It also permits \mention{almost}/\mention{nearly} and occurs after genitives in \mention{her every move}. \textcite[156--158]{spinillo2004reconceptualising} recognizes these differences from \mention{the} and \mention{a}, as well as the articles' greater phonological dependence. Her restricted grouping is thus a substantive alternative, not merely an omission of inconvenient properties.

The issue is how much independent support its shared restrictions provide. Absence of bare, predicative, and partitive uses all limits occurrence without a following nominal; they aren't three independent reasons for a primary category. Little descriptive content extends well beyond the trio. Meanwhile, \mention{a}'s quantificational constructions and \mention{every}'s modifier permissions cross the proposed article boundary. I retain the restricted members within determinative because these connections preserve a broader system without granting unrestricted use to any member.

\textcite[140--144]{spinillo2004reconceptualising} draws an analogy with verbs and prepositions whose complements can be omitted. For demonstratives and quantifiers such as \mention{some} and \mention{all}, occurrence before a nominal (\mention{some teachers}) alternates with occurrence without one (\mention{some left}). The analogy supports category continuity, but leaves the category to be established: \textit{CGEL}, Hudson, and the D-noun analysis can all preserve it.

The argument for retaining the articles is synchronic. Grammaticalization supplies background: \textcite[331--336]{lyons1999} discusses the development of definite articles from demonstratives and of articles expressing cardinality from numerals. These histories can explain changes in content and use, but don't establish the synchronic categorization.

\subsection{Dependent and independent forms}\label{sec:form-selection}

\mention{No}/\mention{none} illustrates a restriction on forms within a paradigm. \textit{CGEL} treats them as inflectional forms of one determinative: dependent \mention{no students} contrasts with independent \mention{none}, while both permit \mention{almost} \citep[389--390]{huddleston2002}. The partitive \mention{none of the students} requires the independent form; \ungram{\mention{no of the students}} is excluded. Form selection remains necessary, just as with \mention{my}/\mention{mine}.

The broader grouping preserves two distinctions: between members' permitted constructions, and between the forms selected within a paradigm. A successful structural account has to respect both. The following section compares the Head relations once an expression's use is licensed.

\section{Ordinary Head, fusion, and modifier attachment}\label{sec:evidence}

The D-noun package pairs the proposed categorization with ordinary Head. This section compares that treatment with fusion in independent \mention{some}, \mention{some of the wine}, and \mention{the lucky few}. Ordinary headedness gives the determinative Head alone; fusion combines Head with a dependent function. The question is how each treatment fits the nominal grammar and preserves the relevant restrictions.

\subsection{The basic Head relations}\label{sec:head-relations}

A \term{nominal} (Nom) contains a head and its internal dependents, excluding an external determiner. In \mention{some apples}, the word \mention{apples} heads Nom, and that Nom heads NP. The same arrangement permits an internal modifier, as in \mention{some red apples}, without making the determining phrase part of the Nom.

Under D-noun, a determinative can head this Nom--NP structure too. Figure~\ref{fig:some} compares the accounts for \mention{take some apples} and \mention{take some}. The dependent expression still has \mention{apples} as its ultimate head. Independent \mention{some} has ordinary Head under D-noun; in \textit{CGEL}, its DP jointly fills Det of NP and Head of Nom.

\begin{figure}[H]
\centering
\textit{CGEL}\par\smallskip
\begin{minipage}{.60\linewidth}\centering
Dependent \mention{some}\par\smallskip
\begin{forest} nominal tree
[VP
 [{\synnode{Head}{V}}, head edge [\mention{take}]]
 [{\synnode{Obj}{NP}}
  [{\synnode{Det}{DP}}
   [{\synnode{Head}{D}}, head edge [\mention{some}]]]
  [{\synnode{Head}{Nom}}, head edge
   [{\synnode{Head}{N\textsubscript{common}}}, head edge [\mention{apples}]]]]]
\end{forest}
\end{minipage}\hfill
\begin{minipage}{.38\linewidth}\centering
Independent \mention{some}\par\smallskip
\begin{forest} nominal tree
[VP
 [{\synnode{Head}{V}}, head edge [\mention{take}]]
 [{\synnode{Obj}{NP}}
  [{\synnode{Head}{Nom}}, head edge, before drawing tree={x+=1.5em}
   [{\synnode{Det--Head}{DP}}, no edge
    [{\synnode{Head}{D}}, head edge [\mention{some}]]]
   {\draw[-] (!uu.south) -- (); \draw[-,line width=1pt] (!u.south) -- (); }]]]
\end{forest}
\end{minipage}

\medskip
\textit{D-noun analysis}\par\smallskip
\begin{minipage}{.60\linewidth}\centering
Dependent \mention{some}\par\smallskip
\begin{forest} nominal tree
[VP
 [{\synnode{Head}{V}}, head edge [\mention{take}]]
 [{\synnode{Obj}{NP}}
  [{\synnode{Det}{NP}}
   [{\synnode{Head}{Nom}}, head edge
    [{\synnode{Head}{N\textsubscript{D}}}, head edge [\mention{some}]]]]
  [{\synnode{Head}{Nom}}, head edge
   [{\synnode{Head}{N\textsubscript{common}}}, head edge [\mention{apples}]]]]]
\end{forest}
\end{minipage}\hfill
\begin{minipage}{.38\linewidth}\centering
Independent \mention{some}\par\smallskip
\begin{forest} nominal tree
[VP
 [{\synnode{Head}{V}}, head edge [\mention{take}]]
 [{\synnode{Obj}{NP}}
  [{\synnode{Head}{Nom}}, head edge
   [{\synnode{Head}{N\textsubscript{D}}}, head edge [\mention{some}]]]]]
\end{forest}
\end{minipage}
\caption{\mention{Take some apples} (left) and \mention{take some} (right) under \textit{CGEL} (top) and the D-noun analysis (bottom). In the upper-right diagram, DP fills Det of NP and Head of Nom. In the lower pair, \mention{some} heads an NP through the same nominal projection; that NP functions as Det on the left and Obj on the right. Function labels appear above category labels. Det marks determiner, Obj object, and Det--Head fusion; the two links show the combined functions. N, D, and V mark noun, determinative, and verb; NP, DP, and VP mark their phrases. Subscripts identify noun subcategories.}\label{fig:some}
\end{figure}

Under D-noun, \term{inheritance} makes Noun's projection rules available to determinatives, subject to their restrictions. Ordinary headedness lets independent determinatives use that nominal structure directly. The separate-D ordinary-Head account can use the same phrase structure by admitting both Noun and D into those rules. The comparison concerns the fit between the lexical grouping and the shared structure.

The genitive NP \mention{Kim's} already fills Det in \mention{Kim's preferences}, as Figure~\ref{fig:genitive} shows. Ordinary determinative headedness gives \mention{some} that phrase type in Det and object uses alike. Plain determinative-headed NPs, such as \mention{almost ten} in \mention{almost ten apples}, then join genitives such as \mention{Kim's} and \mention{my}. The fragment compares the remaining selectional conditions (§\ref{sec:det-uniform}).

\begin{figure}[H]
\centering
\begin{forest} nominal tree
[NP
 [{\synnode{Det}{NP[gen]}}
  [{\synnode{Head}{Nom}}, head edge
   [{\synnode{Head}{N\textsubscript{proper}}}, head edge [\mention{Kim's}]]]]
 [{\synnode{Head}{Nom}}, head edge
  [{\synnode{Head}{N\textsubscript{common}}}, head edge [\mention{preferences}]]]]
\end{forest}
\caption{The genitive NP \mention{Kim's} functions as determiner within \mention{Kim's preferences}, ultimately headed by \mention{preferences}. The representation abstracts from the internal realization of genitive marking.}\label{fig:genitive}
\end{figure}

\subsection{Syntactic completeness and contextual interpretation}\label{sec:saturation}

With a group of people under discussion, \mention{Some left}, \mention{Many came}, and \mention{All agree} illustrate \term{structural saturation}: an argument expression can be syntactically complete without another overt head or determiner. Syntactic completeness differs from contextual interpretation. In \mention{I'll take some}, the relevant substance or set may be supplied by discourse or the situation. Ordinary pronouns also depend on context; both ordinary and fused Head can accommodate that interpretive dependence.

In the following attestation from the CGELBank treebank, \mention{two different Honda models} supplies the domain for the independent object \mention{both}:\footnote{CGELBank \citep{reynolds2023unified}, sentence \href{https://github.com/nert-nlp/cgel/blob/d0a2c2d8c522b23d5aa879fc9cb94948f4285b2c/datasets/ewt.cgel\#L651-L654}{\nolinkurl{reviews-083459-0002}}.}

\ea\label{ex:attested-both}
\mention{Went there yesterday: we are trying to decide between two different Honda models, so we wanted to test-drive both back to back.}
\z

Generalizing expressions such as \mention{Many are called, few are chosen} and \mention{Enough is enough} need no previously uttered common-noun phrase. Interpretation can instead depend on the situation or a generic restriction. Generic pronoun \mention{one}, as in \mention{One shouldn't judge}, provides a parallel without a required overt antecedent.

The comparison with ordinary pronouns and proper nouns is structural as well as interpretive. In \mention{She left} and \mention{Kim left}, \textit{CGEL} permits an NP headed by a nominal with no determiner. Applying that structure to \mention{Some left} preserves the same division between a complete NP and its context-dependent reference. The dependent use of \mention{some} doesn't by itself require a Det function inside every independent occurrence.

Fusion offers another representation of that complete expression. In \textcite{Payne2007}, one constituent jointly realizes Head and a dependent function, subject to structural and adjacency conditions. Contextual recovery alone doesn't decide between this arrangement and ordinary Head. Both analyses retain an overt determinative and can express the restriction supplied by context.

\subsection{Partitives and externally determined expressions}\label{sec:modification}

In \mention{some of the wine}, \mention{some} specifies a quantity drawn from an identifiable amount of wine. The NP \mention{the wine} denotes the partitive domain and is complement of \mention{of}. A pronoun-headed NP or independent genitive NP can express the domain too: \mention{many of them}, referring to previously mentioned people, or \mention{some of Kim's}, referring to Kim's apples.

Compare where the common noun occurs in \mention{some apples} and \mention{some of the wine}. In the partitive, \mention{wine} is embedded inside the \mention{of}-phrase, so it can't be the lexical head of the whole NP. The comparison instead concerns the head-like role of \mention{some} in that larger expression.

Figure~\ref{fig:partitive} gives the D-noun analysis, with the \mention{of}-phrase functioning as complement within Nom, as in \textit{CGEL}'s partitive tree \citep[411--412]{huddleston2002}. \textit{CGEL} represents the head-like role of \mention{some} through Det--Head fusion. Partitives support that role while leaving the choice between fusion and ordinary headedness open.

\begin{figure}[H]
\centering
\begin{forest} nominal tree
[NP
 [{\synnode{Head}{Nom}}, head edge
  [{\synnode{Head}{N\textsubscript{D}}}, head edge [\mention{some}]]
  [{\synnode{Comp}{PP}}
   [{\synnode{Head}{P}}, head edge [\mention{of}]]
   [{\synnode{Comp}{NP}} [\mention{the wine}, roof]]]]]
\end{forest}
\caption{\mention{Some of the wine} under the D-noun analysis. The Head relations lead from the outer NP to \mention{some}. The common noun \mention{wine} is inside the complement PP and doesn't head the whole expression. The inner NP is abbreviated.}\label{fig:partitive}
\end{figure}

Both Head analyses permit the overt quantifier and domain to supply a compositional interpretation of \mention{some of the wine}.

Independent \mention{few} and certain other indefinite determinatives permit definite determination. Under D-noun, \mention{few} heads an ordinary NP both alone and in \mention{the few}, where \mention{the} fills Det. Adding an optional adjective gives \mention{the lucky few}, with the same Head and dependent functions as \mention{the lucky survivors}.

\textit{CGEL} explicitly permits determinatives used as internal modifiers to fuse with Head, as in \mention{the other two} and \mention{these few here} \citep[415--416]{huddleston2002}. Its analysis of \mention{the few mistakes} assigns \mention{the} to Det and \mention{few} to Mod \citep[392]{huddleston2002}. That dependent use supplies the counterpart for a Mod--Head analysis of independent \mention{few} after an external determiner.

Figure~\ref{fig:few} compares the resulting analyses of \mention{the lucky few}. Both have an overt Det and a Nom modified by \mention{lucky}. In the D-noun analysis, \mention{few} fills Head. In the fusion analysis, its DP fills Mod--Head. The competing structure thus uses machinery already needed for adjectival fusion; it requires neither a silent noun nor conversion.

\begin{figure}[H]
\centering
\begin{minipage}{.48\linewidth}\centering
\begin{forest} nominal tree
[NP
 [{\synnode{Det}{NP\textsubscript{D}}} [\mention{the}, roof]]
 [{\synnode{Head}{Nom}}, head edge
  [{\synnode{Mod}{AdjP}} [\mention{lucky}]]
  [{\synnode{Head}{N\textsubscript{D}}}, head edge [\mention{few}]]]]
\end{forest}
\end{minipage}\hfill
\begin{minipage}{.48\linewidth}\centering
\begin{forest} nominal tree
[NP
 [{\synnode{Det}{DP}} [\mention{the}, roof]]
 [{\synnode{Head}{Nom}}, head edge
  [{\synnode{Mod}{AdjP}} [\mention{lucky}]]
  [{\synnode{Head}{Nom}}, head edge, before drawing tree={x+=1.5em}
   [{\synnode{Mod--Head}{DP}}, no edge
    [{\synnode{Head}{D}}, head edge [\mention{few}]]]
   {\draw[-] (!uu.south) -- (); \draw[-,line width=1pt] (!u.south) -- (); }]]]
\end{forest}
\end{minipage}
\caption{\mention{The lucky few} with ordinary Head (left) and Mod--Head fusion (right). The determiner is separately realized in both. Internal structure within the article phrase and AdjP is suppressed.}\label{fig:few}
\end{figure}

Table~\ref{tab:matched-heads} compares the Head relations across five constructions; §\ref{sec:det-uniform} supplies the rules under both taxonomies.

\begin{table}[H]
\centering\small
\caption{Head and dependent functions under matched analyses. Each independent expression is an NP; \mention{the} fills Det wherever it appears. The ordinary-Head column applies to both taxonomies permitting ordinary determinative heads.}\label{tab:matched-heads}
\begin{tabular}{>{\raggedright\arraybackslash}p{2.8cm}>{\raggedright\arraybackslash}p{4.5cm}>{\raggedright\arraybackslash}p{4.5cm}}
\toprule
Expression & Separate D with fusion & Ordinary determinative Head \\
\midrule
\mention{few survivors} & \mention{few}: Det; \mention{survivors}: Head & Same functions \\
Independent \mention{few} & \mention{few}: Det--Head & \mention{few}: Head \\
\mention{the few} & \mention{few}: Mod--Head & \mention{few}: Head \\
\mention{the lucky few} & \mention{few}: Mod--Head; \mention{lucky}: Mod & \mention{few}: Head; \mention{lucky}: Mod \\
\mention{the idle rich} & \mention{rich}: Mod--Head; \mention{idle}: Mod & Same fusion and modifier functions \\
\bottomrule
\end{tabular}
\end{table}

Fusion groups \mention{the lucky few} with \mention{the idle rich}; ordinary headedness groups it with \mention{the lucky survivors} and unifies bare and externally determined \mention{few}. Both cover the example. The choice is between sharing fusion across D and adjective, and sharing ordinary Head structure across determinative constructions.

Reducing fusion's applications can simplify this description even when fusion remains available for adjectives. The advantage is the uniform treatment of \mention{few}; its value depends on any additional conditions or lost generalizations elsewhere.\footnote{Compare the theory of second best in \textcite[11--12]{lipseylancaster1956secondbest}: under a constraint preventing an optimum, satisfying more optimality conditions needn't improve the outcome. The analogy concerns interactions among grammatical choices, not a formal optimum for the grammar.} Section~\ref{sec:costs} compares those costs while holding the lexical restrictions fixed.

The constructed \mention{the remaining three} extends the modifier pattern to cardinals. It doesn't independently decide whether \mention{three} is determinative or has a common-noun use: \textcite{reynolds2026numerals} argues for both uses of cardinals. Section~\ref{sec:costs} treats their unification as a consequence conditional on that analysis.

The modifier permissions needed under ordinary headedness now require a closer comparison.

\subsection{Modifier selection and attachment}\label{sec:payne}

The selection and ordering contrast in §\ref{sec:adjectival-profile} guides the attachment analysis. Approximative and focusing modifiers have established NP parallels; the degree series selects the four gradable quantifiers. Both ordinary-Head accounts can represent that difference by distinguishing NP-peripheral modification from modification inside Nom.

Peripheral modifiers attach outside an NP's internal dependents. Compare \mention{only you}, \mention{for almost my entire life}, and \mention{Usually a careful driver, Anne found her mind wandering}. The modifiers precede any determiner belonging to the NP they modify. \textit{CGEL} already assigns them NP-level attachment \citep[430--431]{huddleston2002}. Focusing \mention{even} and \mention{just}, and approximatives such as \mention{nearly}, \mention{hardly}, \mention{virtually}, and \mention{practically}, likewise have NP uses, subject to their own scope and selectional restrictions.

In \mention{almost every experienced teacher}, \mention{almost} belongs with \mention{every}, while \mention{experienced} modifies \mention{teacher}. Figure~\ref{fig:every} places \mention{almost} at the periphery of the NP headed by \mention{every}. That smaller NP functions as Det in the larger NP, whose ultimate head is \mention{teacher}.

\begin{figure}[H]
\centering
\begin{forest} nominal tree
[NP
 [{\synnode{Det}{NP\textsubscript{D}}}
  [{\synnode{Mod}{AdvP}} [\mention{almost}]]
  [{\synnode{Head}{NP\textsubscript{D}}}, head edge
   [{\synnode{Head}{Nom}}, head edge
    [{\synnode{Head}{N\textsubscript{D}}}, head edge [\mention{every}]]]]]
 [{\synnode{Head}{Nom}}, head edge
  [{\synnode{Mod}{AdjP}} [\mention{experienced}]]
  [{\synnode{Head}{N\textsubscript{common}}}, head edge [\mention{teacher}]]]]
\end{forest}
\caption{Peripheral \mention{almost} modifies the smaller NP headed by \mention{every}; the resulting \mention{almost every} NP functions as Det. The adjective \mention{experienced} modifies the common-noun nominal. Internal structure within the one-word modifier phrases is suppressed.}\label{fig:every}
\end{figure}

\textcite[40--42]{payne2010} defend keeping \mention{few}, \mention{any}, and related forms in one lexical category across dependent and independent uses. In \mention{hardly any money} and independent \mention{hardly any}, \mention{hardly} remains an adverb and \mention{any} retains its category. The D-noun analysis preserves that continuity. It also preserves their contrast between adverbial premodification, as in \mention{almost anybody}, and adjectival postmodification, as in \mention{nothing absolute}.

The proposed analysis gives \mention{very} internal attachment in dependent and independent uses alike. In \mention{the very few people}, it modifies the Nom headed by \mention{few}; the resulting \mention{very few} NP functions as Mod in the larger NP. In \mention{the very few who objected}, \mention{few} remains ordinary Head and \mention{very} remains internal. The degree relation is unchanged, and one internal permission covers both constructions. The same analysis applies to \mention{the very many}.

NP degree modifiers provide a further connection. Determinatives permit determinative premodifiers such as \mention{this} in \mention{this much} \citep[393]{huddleston2002}. Comparative determinatives also take established NP modifiers: \mention{a lot} in \mention{a lot fewer}, alongside \textit{CGEL}'s \mention{a lot more than fifty} \citep[432]{huddleston2002}. These modify the quantity expressed by the head. Under the proposed analysis, they're internal NP modifiers within its nominal projection, with permissions specific to the head and construction.

\textcite[42--47]{payne2010} establish that adverbs can postmodify common nouns, as in the constructed \mention{the changes globally to the climate}. They also analyse \mention{almost} as modifying the attributive nominal \mention{textbook} in \mention{an almost textbook case} \citep[p.~75, n.~3]{payne2010}. Adverbial modification therefore isn't categorically excluded from Noun, even internally. Its broader availability with determinatives remains a difference of distribution within the proposed category.

Approximatives need internal attachment where they follow an external determiner on an independent cardinal: \mention{the almost thirty who came}. The internal AdvP permissions at issue are thus the degree series on \mention{few}, \mention{many}, \mention{much}, and \mention{little} across dependent and independent uses, and approximatives on externally determined independent cardinals.

\subsection{Compounds and their modifier domains}\label{sec:compounds}

A more demanding modifier comparison is \mention{hardly anyone present}. \textcite[581--583]{Payne2007} assign \mention{hardly} to DP structure and \mention{present} to nominal structure. The compound \mention{anyone} takes the premodifiers of its determinative base \mention{any}: compare \mention{hardly any writer present}. The adjective realizes a specialized \term{restrictor} function, restricted to post-head position and non-recursive. Figure~\ref{fig:compound} contrasts this account with an ordinary-Head analysis.

\begin{figure}[H]
\centering
\begin{minipage}{.48\linewidth}\centering
\begin{forest} nominal tree
[NP
 [{\synnode{Mod}{AdvP}} [\mention{hardly}]]
 [{\synnode{Head}{NP}}, head edge
  [{\synnode{Head}{Nom}}, head edge
   [{\synnode{Head}{N\textsubscript{D}}}, head edge [\mention{anyone}]]
   [{\synnode{Mod}{AdjP}} [\mention{present}]]]]]
\end{forest}
\end{minipage}\hfill
\begin{minipage}{.48\linewidth}\centering
\begin{forest} nominal tree
[NP
 [{\synnode{Head}{Nom}}, head edge, before drawing tree={x+=4em}
  [{\synnode{Det--Head}{DP}}, no edge
   [{\synnode{Mod}{AdvP}} [\mention{hardly}]]
   [{\synnode{Head}{D}}, head edge [\mention{anyone}]]]
  {\draw[-] (!uu.south) -- (); \draw[-,line width=1pt] (!u.south) -- (); }
  [{\synnode{Mod}{AdjP}} [\mention{present}]]]]
\end{forest}
\end{minipage}
\caption{\mention{Hardly anyone present} with ordinary Head (left) and fusion (right). On the left, \mention{hardly} is peripheral to NP and \mention{present} is internal to Nom. Both post-head Mod positions realize the specialized restrictor. In the fusion tree, DP fills Det of NP and Head of Nom, following \textcite[p.~582, (13d)]{Payne2007}. Modifier-phrase interiors are suppressed.}\label{fig:compound}
\end{figure}

In the ordinary-Head tree, \mention{anyone} inherits nominal projection from Noun and its premodifier permissions from its determinative base. Peripheral \mention{hardly} follows the approximative pattern and its NP parallels in §\ref{sec:payne}. The compound construction independently excludes external determination and supplies the internal post-head restrictor permission, with the same position and non-recursion conditions as the fusion account. It doesn't license a corresponding pre-head adjective.

The relative clauses in §\ref{sec:independent} are ordinary postmodifiers, distinct from the specialized restrictor. When both occur, the relative follows it: \mention{something useful that I found}. \textit{CGEL} explicitly gives compounds the common noun's range of ordinary postmodifiers while preserving this ordering condition \citep[423]{huddleston2002}.

The same division applies to \mention{someone} and \mention{everybody}: the compound construction licenses the post-head restrictor, while the determinative base supplies its premodifier permissions. The attested \mention{something reliable and good looking} illustrates the post-head pattern (§\ref{sec:independent}).

Both accounts separate the modifier domains structurally: fusion uses the DP--Nom boundary, and ordinary headedness uses the NP--Nom boundary. Their constituent groupings differ: ordinary Head groups \mention{anyone present} in the inner NP, whereas fusion groups \mention{hardly anyone} in DP.

\subsection{Independent genitives}\label{sec:genitive-head}

\textit{CGEL} also uses fusion with nouns. Anaphoric \mention{mine}, understood as \enquote*{my car}, combines the possessive relation with an understood nominal description and receives a fused analysis. Predicative \mention{mine}, as in \mention{it's mine}, can instead mean \enquote*{belongs to me}: it expresses the relation directly, without an understood nominal head \citep[410--411, 470--471]{huddleston2002}. A suitable context may permit either reading.

The proposed analysis gives \mention{mine} ordinary Head in both anaphoric and predicative uses. Its genitive form expresses a relation to the speaker; context can supply the understood nominal description, as it can with independent \mention{few}. The \mention{my}/\mention{mine} alternation remains a form-selection condition, parallel to \mention{no}/\mention{none}. This extends ordinary headedness to an existing noun subcategory.

In referential uses, the independent-genitive construction supplies the NP's interpretation and agreement properties based on the possessed entity or entities. Compare constructed \mention{I am ready} with \mention{Mine is ready} (\enquote*{my contribution}) and \mention{Mine are ready} (\enquote*{my slides}). The pronoun's first-person feature identifies the possessor; the independent NP has third-person singular or plural agreement.

The fragment now makes the permissions and restrictions explicit so that the structural accounts can be compared.

\section{The matched fragment and comparative costs}\label{sec:economy}

A grammatical \term{fragment} states rules for a specified range of constructions. The four accounts receive the same judgments and lexical restrictions. Comparing the rules needed for coverage makes their commitments explicit; it doesn't test them against new data.

\subsection{Scope and use permissions}\label{sec:fragment}

The fragment covers determination, ordinary subject and object uses, partitives, and the modifier and compound comparisons in §\ref{sec:evidence}. It admits relative clauses as postmodifiers but doesn't analyse their internal grammar. Predication, interrogative clause structure, and modification outside NP structure lie beyond its scope. Degree uses enter the wider accounting in §\ref{sec:enough}.

Table~\ref{tab:permissions} records \term{use permissions}. \term{Dependent} means that the form heads a phrase in Det function before another nominal. \term{Independent} means that its phrase occurs without another nominal in the displayed subject or object uses, as in \mention{Some left} and \mention{I saw some}. These labels describe the constructions covered here, rather than every possible function. Quotation, metalinguistic naming, and subordinate-clause uses of dependent genitives fall outside the fragment. Noun membership doesn't itself grant a use permission.

\begin{table}[H]
\centering\small
\caption{Shared permissions in the fragment. For \mention{some}, the target column concerns its unstressed use before a noun. The inventory is deliberately limited to the displayed constructions.}\label{tab:permissions}
\begin{tabular}{lcc>{\raggedright\arraybackslash}p{6.2cm}}
\toprule
Form & Dependent & Independent & Target in the Det construction \\
\midrule
\mention{the} & yes & no & Singular or plural; count or non-count \\
\mention{a}, \mention{every} & yes & no & Singular count nominal \\
\mention{some} & yes & yes & Plural count or non-count nominal \\
\mention{few} & yes & yes & Plural count nominal \\
\mention{my} & yes & no & No count or number restriction \\
\mention{she} & no & yes & Not applicable \\
\bottomrule
\end{tabular}
\end{table}

Further constructional restrictions apply: \mention{she} is a subject form, whereas the corresponding ordinary object form is \mention{her}. The target restrictions concern the common-noun nominal being determined: \mention{a} requires a singular count target such as \mention{book} in \mention{a book}.

Hudson distinguishes dependency from external headedness. He argues that determiner and common noun depend on each other, although only one connects the phrase to its surroundings \citep[7--8]{hudson2004determiners}. His temporal adjunct evidence supports common-noun headedness: \mention{I saw him that day} is possible, whereas \ungram{\mention{I saw him that point in time}} isn't, despite the similar temporal meanings \citep[10--12]{hudson2004determiners}. The lexical noun matters, as well as the construction's determiner restrictions.

The reciprocal selection facts are substantial: a singular count noun normally requires determination, while \mention{every} requires an overt nominal and \mention{some} doesn't. They establish conditions in both directions. The noun can remain Head while both constituents impose these conditions.

\subsection{Shared phrase types and projection rules}\label{sec:det-uniform}

In \mention{some apples} and \mention{Kim's apples}, the determining phrases share the NP category under the proposed package. Separate D can permit the same NP projection. For the D-noun/fusion alternative, I retain that projection and add fusion in independent uses. These three implementations therefore share an inventory of phrase types. \textit{CGEL} instead distinguishes DP from genitive NP.

\textit{CGEL} also admits a restricted range of plain-case NPs and PPs as determiners, as in \mention{what size hat} and \mention{over thirty ties} \citep[Ch.~5, §4]{huddleston2002}. The phrase types are:

\[
\begin{aligned}
\text{\textit{CGEL}:} &\quad \mathrm{Det}:\{\mathrm{DP},\mathrm{NP},\mathrm{PP}\} \\
\text{The three NP-projecting accounts:} &\quad \mathrm{Det}:\{\mathrm{NP},\mathrm{PP}\}
\end{aligned}
\]

Determinative phrases and genitive NPs supply the principal realizations in \textit{CGEL}; PPs are restricted. Giving determinatives NP projection consolidates those principal types. Determinative-headed, genitive, and other licensed NPs still require separate identification. \textit{CGEL} already states definiteness and the single-Det restriction functionally, so recategorization doesn't derive them. The following fragment compares Head relations with those restrictions held fixed.

The four combinations permit comparison of the Head analyses within each taxonomy and of the taxonomies under each Head analysis. Table~\ref{tab:economy} summarizes their commitments. The implementations below specify the phrase types as well as the lexical categories and Head functions.

The D-noun analysis applies nominal projection to determinatives through their membership in Noun. A separate condition checks use permissions. In the schemata below, $h$ identifies the lexical head throughout its projection; parentheses mark optional dependents. The head's entry and construction restrict every dependent. $\mathrm{Mod}_{\mathrm{periph}}$ adds the peripheral premodifier discussed in §\ref{sec:payne}; the compact rule leaves further NP layering implicit.

$\mathrm{Mods}_{\mathrm{post}}$ denotes a possibly empty, constructionally ordered sequence of postmodifiers. For compounds, this permits one specialized restrictor followed by ordinary postmodifiers: \mention{useful} precedes \mention{that I found} in \mention{something useful that I found}. It doesn't permit unrestricted repetition of the restrictor. The same sequence notation and ordering conditions apply under all four accounts.

\[
\begin{aligned}
\mathrm{Nom}_{h} &\to (\mathrm{Mod}_{\mathrm{pre}})\quad \mathrm{Head}:\mathrm{N}_{h}\quad (\mathrm{Comp})\quad \mathrm{Mods}_{\mathrm{post}} \\
\mathrm{NP}_{h} &\to (\mathrm{Mod}_{\mathrm{periph}})\quad (\mathrm{Det})\quad \mathrm{Head}:\mathrm{Nom}_{h} \\
\operatorname{licensed}(\mathrm{NP}_{h},f,c)\quad &\Longleftrightarrow\quad f\in U_h\ \land\ C_h(f,c)
\end{aligned}
\]

Here $U_h$ is the head's set of use permissions, $f$ is the completed NP's external function, and $C_h$ checks the lexical and constructional conditions in context $c$. For Det use, these include compatibility with the target nominal; for a singular count argument, they include required determination. Optionality in the second rule doesn't override those conditions. Intermediate Head links within a projection don't require independent argument permission.

These checks accommodate the reciprocal selection discussed by Hudson. The determining phrase checks the target nominal's count and number properties; the outer noun's projection checks whether determination is required and supplied. Both conditions have to hold even though the tree assigns only one Head to each phrase.

In \mention{some apples}, \mention{some} heads an NP whose Det permission and plural-count target requirement are satisfied. \mention{Apples} heads the outer Nom, which heads the NP. In \mention{Some left}, \mention{some} heads an NP whose argument permission is satisfied. \mention{Every apple} passes the Det and singular-count checks; ordinary independent \ungram{\mention{Every arrived}} fails the argument-permission check. \mention{The apple} and \ungram{\mention{The arrived}} differ in the same way.

Each NP's use permissions follow its own head. In \mention{the apple}, the article's Det permission licenses the dependent NP headed by \mention{the}. The outer NP takes its argument permission from \mention{apple}, whose requirement for determination is satisfied by the article.

The same distinction between phrase levels matters beyond the fragment. In \mention{I know which apple it is}, \mention{which} heads the determining phrase inside \mention{which apple}. The whole \mention{which apple} is the preposed predicative complement of \mention{is} in the embedded interrogative clause. Its function doesn't make \mention{which} independent; \mention{apple} still heads that NP.

A singular count common noun such as \mention{book} faces a different restriction from an article. In \mention{a book}, its requirement for determination is satisfied; bare \ungram{\mention{Book arrived}} leaves that requirement unsatisfied. \mention{Books arrived} has no such requirement. Requiring determination for a common noun doesn't itself block an article-headed NP from argument use. An article's exclusion from argument use must still be stated separately.

\mention{Some} and \mention{few} permit a partitive \mention{of}-phrase within their nominal projection. Independent \mention{few} permits definite determination in \mention{the few} and optional adjectival modification in \mention{the lucky few}. Degree \mention{very} modifies \mention{few} internally across dependent and independent uses.

Peripheral \mention{almost} modifies the NP headed by \mention{every}; this doesn't license \mention{experienced} as its modifier. The cardinal construction \mention{the almost thirty who came} licenses internal approximative \mention{almost}. In \mention{hardly anyone present}, \mention{hardly} is peripheral to NP and the restrictor \mention{present} is internal to Nom.

Under separate D with ordinary Head, the NP and use-permission rules are the same. Its Nom rule admits either Noun or determinative as lexical head:

\[
\mathrm{Nom}_{h}\to(\mathrm{Mod}_{\mathrm{pre}})\quad\mathrm{Head}:\{\mathrm{N}_{h},\mathrm{D}_{h}\}\quad(\mathrm{Comp})\quad\mathrm{Mods}_{\mathrm{post}}
\]

This projection applies to dependent as well as independent determinatives, eliminating DP from the fragment while retaining primary D. Determiner selection still identifies D-headed NPs and suitable genitives. In \mention{the lucky few}, \mention{few} is D in ordinary Head function. In \mention{hardly anyone present}, the determinative base and compound construction supply the same modifier restrictions as the D-noun account. Adjectival Mod--Head fusion remains available in both.

The separate-D fusion account relates independent determinatives to their dependent counterparts. In addition to nominal projection for nouns, it uses DP projection and permits a DP to realize a fused function in nominal structure. The partitive complement and compound restrictor belong to Nom, following the trees in \textit{CGEL} and \textcite[582]{Payne2007}:

\[
\begin{aligned}
\mathrm{DP}_{h} &\to (\mathrm{Mod}_{\mathrm{pre}})\quad \mathrm{Head}:\mathrm{D}_{h}\quad (\mathrm{Comp}) \\
\mathrm{Nom}_{h} &\to (\mathrm{Mod}_{\mathrm{pre}})\quad F:\mathrm{DP}_{h}\quad (\mathrm{Comp})\quad \mathrm{Mods}_{\mathrm{post}} \\
F&\in\{\mathrm{Det\text{--}Head},\mathrm{Mod\text{--}Head}\}
\end{aligned}
\]

The ordinary NP rule embeds this Nom. Det--Head jointly realizes Det of NP and Head of Nom, excluding a second Det; Mod--Head fills an internal modifier's function and Head of Nom, allowing external determination. Lexical and constructional conditions select the fused function and permitted dependents. These schemata cover independent \mention{few}, \mention{few of them}, \mention{the lucky few}, and \mention{hardly anyone present}.

The D-noun/fusion alternative uses the NP already projected by a determinative noun in place of DP in the fusion rule. That NP jointly fills Det of an outer NP and Head of its Nom, or Mod--Head where external determination is present. Fusion applies once to the determinative's projection. The partitive complement and compound restrictor remain on the outer Nom, as in the separate-D fusion account.

This version of fusion adds nominal structure around a determinative expression that already forms an NP. Ordinary headedness uses that NP directly, giving bare, partitive, and externally determined determinatives the familiar nominal Head relations.

Both ordinary-Head accounts provide the same structures. Under D-noun, those structures instantiate the usual projection of a member of Noun; separate D extends that projection across a primary-category boundary. I take this alignment of category and structure to favour the D-noun package. Their shared coverage leaves both accounts viable, but doesn't make their categorizations equally well supported.

A grammar with Hudson's nested categorization can use the same permissions and projection rules. Holding those rules fixed, placing the broader pronoun category on determinative's inheritance path from Noun leaves the fragment's judgments unchanged.

\subsection{Degree uses outside noun phrases}\label{sec:enough}

Determinatives also modify expressions outside NP structure. The three NP-projecting implementations retain that projection in degree uses. The relevant comparison is with the degree uses of established NPs, whose distribution already extends beyond argument and determiner functions.

\mention{Enough} tests both position and external function. It precedes a nominal in \mention{enough money} and can follow one in \mention{money enough}; post-head \mention{enough} can't itself be premodified, as shown by \ungram{\mention{money almost enough}} \citep[396--397, 445]{huddleston2002}. Only the permitted construction licenses each position.

\mention{Enough} also modifies adjectives, adverbs, verbs, and some PPs: \mention{good enough}, \mention{quickly enough}, \mention{I hadn't prepared enough}, and \mention{enough in control} \citep[396--397]{huddleston2002}. The degree determinatives \mention{much} and \mention{little}, and \mention{no}/\mention{none}, likewise have uses outside NP structure \citep[390, 395--397]{huddleston2002}. These permissions must survive recategorization. Noun membership doesn't confine every projection to argument or determiner function.\footnote{\mention{Both}, \mention{either}, and \mention{neither} also serve as markers of coordination \citep[1305, 1308]{huddleston2002}. Their noun categorization would preserve that further category--function combination; the present fragment doesn't analyse coordination.}

NPs already modify adjectives in \mention{three years old}, \mention{a great deal smaller}, and \mention{plenty big enough}. \textit{CGEL} also explicitly treats \mention{lots better} and \mention{heaps worse} as containing quantificational NP modifiers \citep[549--550]{huddleston2002}. These uses extend the connection with the quantificational common nouns compared in §\ref{sec:quant-nouns}.

Attributive uses are attested too: \mention{some heaps better photo's} and \mention{some lots better ones}.\footnote{\href{https://www.aulro.com/afvb/motor-cycling/61858-my-track-day-phillip-island-motogp-track-post798794.html}{AULRO, moose, 19 August 2008, post 2}; original spelling retained. The \href{https://www.uberpeople.net/threads/let-s-play-would-you-have-taken-this-trip.395130/}{Uber Drivers Forum} example was verified in an indexed excerpt; its author and post date remain unverified.} In \textit{CGEL}'s analysis, \mention{ones} here is an anaphoric common noun. Both examples therefore have attributive \mention{better} with a nominal degree modifier. Measure \mention{miles} supplies parallels in \mention{some miles better front brakes} and \mention{some miles-better bands}.\footnote{\href{https://www.uksaabs.co.uk/UKS/viewtopic.php?t=131481\#p1256815}{UKSaabs, rallyv4, 23 January 2014, post 6}; Shaun Ryder in \href{https://www.newstatesman.com/culture/2019/04/i-look-like-uncle-fester-the-second-life-of-shaun-ryder}{Kate Mossman's interview, \textit{New Statesman}, 17 April 2019}. The latter retains the source's hyphen, which doesn't itself settle the modifier's syntactic structure.}

These attestations challenge \textit{CGEL}'s general exclusion of NP modifiers from attributive AdjPs. Its contrast between predicative \mention{a great deal better} and \ungram{\mention{some a great deal better proposals}}, against permitted \mention{some much better proposals}, remains a more restricted observation \citep[551--552]{huddleston2002}. The excluded example also juxtaposes \mention{some} and \mention{a}, so its unacceptability can't by itself be attributed to NP category.

\textit{CGEL} already permits \mention{She's a lot better player than me}, analysing the article belonging to \mention{a lot} as lost after the article determining the larger NP \citep[p.~552, n.~8]{huddleston2002}. The \mention{heaps} and \mention{lots} examples have outer \mention{some} and no internal article, so they extend beyond this exception.

The attestations don't establish unrestricted attributive use for nominal degree modifiers. They require a closer comparison of lexical, constructional, and register conditions under all four implementations. NP projection places degree determinatives within an existing nominal range; the attributive contrast doesn't yet establish an added cost, or a corresponding advantage for retaining DP.

\subsection{The comparative judgment}\label{sec:costs}

Table~\ref{tab:economy} compares the four implementations. The Head configurations, modifier permissions, and taxonomic consequences are kept distinct so that the contribution of each choice is visible. The inventory, constructions, and readings remain fixed.

\begin{table}[H]
\centering\small
\caption{Four implementations compared. All retain the lexical and constructional restrictions and adjectival fusion, and must account for attributive nominal degree modifiers. Cardinal and \mention{one} groupings are consequences of the categorization.}\label{tab:economy}
\setlength{\tabcolsep}{3pt}
\begin{tabular}{>{\raggedright\arraybackslash}p{2.5cm}>{\raggedright\arraybackslash}p{3.6cm}>{\raggedright\arraybackslash}p{3.6cm}>{\raggedright\arraybackslash}p{2.65cm}}
\toprule
Account & Independent determinatives & Modifier permissions & Taxonomic consequences \\
\midrule
D-noun, ordinary Head & Nom--NP with Head alone across bare, partitive, and externally determined uses. & Peripheral NP premodifiers; head-specific internal modifiers. Degree determinatives project NP. & D within Noun. Cardinal uses and the three \mention{one} lexemes are grouped there. \\
Separate D, ordinary Head & Same structures, with a Nom rule admitting either N or D. & Same internal and peripheral permissions and degree projection as the proposed package. & D remains separate from nominal counterparts. \\
Separate D, fused Head & DP fills Det--Head or Mod--Head in nominal structure. & DP--Nom separates premodifiers and postmodifiers. Degree determinatives project DP. & Same primary-category divisions as separate D with ordinary Head. \\
D-noun, fused Head & A projected NP fills Det--Head or Mod--Head in additional nominal structure. & Premodifiers within the fused NP; postmodifiers in the outer Nom. Degree determinatives project NP. & Same Noun grouping as the proposed package. \\
\bottomrule
\end{tabular}
\end{table}

The articles' restrictions, \mention{no}/\mention{none} form selection, and the compound restrictor conditions are already required in the separate-D grammar. Retaining them isn't an added cost of D-noun. Likewise, differences among existing noun subcategories in determination and modification remain in place (Table~\ref{tab:existing}).

The internal degree-AdvP permissions on \mention{few}, \mention{many}, \mention{much}, and \mention{little}, and the internal approximatives with externally determined cardinals, remain explicit (§\ref{sec:payne}). Both ordinary-Head accounts also use peripheral NP modification. Fusion instead places the compound's premodifiers within the fused phrase and its postmodifiers on the outer Nom (§\ref{sec:compounds}).

\textcite{reynolds2026numerals} distinguishes determinative numerals such as \mention{ten} in \mention{ten men}, proper-noun cases such as \mention{10} in \mention{Room 10}, and common-noun numerals such as \mention{tens} in \mention{tens of pens}. Under the D-noun categorization, all fall within primary Noun, retaining their subcategory and constructional differences. Ordinals remain adjectives, and complex numeral phrases remain distinct from single lexemes.

\textcite[797--798]{payne2013anaphoric} distinguish three lexemes spelled \mention{one}: determinative, anaphoric common noun, and generic pronoun. These remain distinct under the D-noun categorization, but all belong within primary Noun.

Complex cardinals also separate category from function. In \mention{two hundred books}, the whole \mention{two hundred} fills Det; internally, \mention{two} modifies the magnitude head \mention{hundred} \citep[§4]{reynolds2026numerals}. In \mention{these two hundred books}, \mention{these} fills Det and \mention{two hundred} is an internal modifier. One Det function doesn't entail a limit of one determinative lexeme per NP.

These cardinals and the lexemes spelled \mention{one} retain their distinctions within one primary category. That consequence belongs to the categorization and is shared by its ordinary- and fused-Head implementations.

I favour the D-noun package because the profile and the ordinary NP structures support a common nominal treatment. Including determinatives within Noun makes those structures a natural application of nominal grammar; their regularity in turn strengthens the grouping. The commitments are complementary, although neither requires the other. Separate D with ordinary Head and D-noun with fusion remain coherent alternatives. Reducing the number of primary categories alone wouldn't establish the preference.

If determinative-headed and genitive expressions require different projection rules after their independently motivated restrictions are held fixed, the shared-projection proposal in §\ref{sec:det-uniform} loses its advantage. That would favour separate phrase types. Retaining a separate primary D requires the further case that the category boundary captures the recurring differences better than a determinative subcategory within Noun.

\section{Coordinate or nested subcategories}\label{sec:inheritance}

Including determinatives within Noun leaves a further question: where within Noun do they belong? In the proposed hierarchy, common noun, proper noun, pronoun, and determinative are four coordinate subcategories. Hudson instead puts determinative inside pronoun, which is itself inside Noun. Both group determinatives with nouns. They differ in whether pronouns and determinatives form an intermediate category that excludes common and proper nouns.

Hudson makes taxonomic inclusion explicit. His noun category contains common noun, proper noun, and pronoun \citep[253--254]{hudson2010wordgrammar}; a determiner is a pronoun whose \term{valency} permits the relevant common-noun dependent \citep[9--10]{hudson2004determiners}. The word's category remains constant across independent and dependent uses, which differ in their permitted dependents.

Hudson's criteria centre on licensing a singular count common noun and on mutual exclusion in that use. He sets aside \mention{all}, cardinal numerals, and quantifiers restricted to plural or non-count nouns \citep[9--10]{hudson2004determiners}. Their place in his taxonomy remains unsettled by those criteria. Extending the pronoun analysis to \textit{CGEL}'s fuller inventory requires comparing its modifier selection, grade, and quantitative meanings with the existing pronouns' properties.

The intermediate category would be useful if it supported grammatical generalizations applying to pronouns and determinatives together. The coordinate analysis can still record properties shared by those two groups; it doesn't claim that every noun subcategory is equally similar to every other. The issue is whether their shared properties warrant another level in the lexical hierarchy.

\textcite{reynolds2021} compares 138 word forms through properties recorded as present or absent, such as accepting \mention{almost}. Its unsupervised clustering groups forms by these similarities without being given their category labels. The groups broadly resemble the pronoun and determinative inventories, though the outcome varies with initialization and feature selection.\footnote{The study reports 232 properties, whereas the public file contains 155. The accompanying \href{run:matrix-audit.pdf}{\textit{Replication audit of the English determinative--pronoun feature matrix}} documents the discrepancy, reproduces the published statistical decomposition, and examines sensitivity. The public file is preserved unchanged.} But the study contains no common or proper nouns. Distinguishing the two groups therefore establishes neither their taxonomic rank nor the proposed superordinate category.

The grammatical case for nesting needs to be assessed on its own. Many pronouns and determinatives have little descriptive content and depend on context for interpretation. Pro-form gender supplies further connections (§\ref{sec:inflection}). These properties make the grouping plausible, but they don't uniquely identify it: primary naming uses of proper nouns also depend on context, and pro-form gender extends to non-nominal expressions \citep{reynolds2025proformgender}.

Restrictions on determination and modification are another candidate. Personal pronouns generally resist external determiners and permit few internal modifiers: \mention{poor old me} illustrates the restricted adjectival pattern \citep[429--430]{huddleston2002}. But proper names also resist free determination and modification in their primary uses \citep[517, 519--520]{huddleston2002}. These restrictions don't pick out pronouns and determinatives alone.

Nor are those restrictions uniform within the proposed intermediate category. Determinatives allow \mention{the few}, \mention{the two}, and \mention{these three}, with item-specific conditions; \textit{CGEL} gives \mention{these few here} and \mention{the many who did} \citep[415--416]{huddleston2002}. Their adverbial modifier patterns differ from the restricted adjectival pattern of personal pronouns. A broader pronoun category couldn't simply pass the personal pronouns' rules down to determinatives.

I therefore place determinative alongside pronoun. Their shared nominal structure belongs at the Noun level, while their more specific patterns remain distinct: person, case, and reflexivity organize much of the pronoun system; quantification, determination, and degree organize much of the determinative system. Nesting remains a coherent alternative, but the comparisons here don't establish a further shared set of grammatical rules that requires it.

Both inheritance paths supply the fragment's structures (§\ref{sec:det-uniform}). The coordinate preference rests on the more specific profile; accepting Noun membership and ordinary Head leaves this further choice open.

\section{Conclusion}\label{sec:conclusion}

I propose including English determinatives within Noun. The strongest positive comparison comes from quantificational common nouns: their selected complements, number transparency, restricted dependents, and degree uses connect with parts of the determinative inventory. The connected adjectival profile of the four gradable quantifiers deserves weight on the same terms. Broader constructional and referential connections, with partial inflectional corroboration, favour the nominal grouping.

The grouping includes restricted members through their integration into the determinative system. Articles needn't independently exhibit the whole profile, and \mention{no}/\mention{none} retains its form-selection conditions. The argument for membership thus concerns a category with internal differences, rather than unrestricted use by every form.

Ordinary Head gives bare, partitive, and externally determined determinatives the familiar Nom--NP structure. That regularity complements the profile case for including them within Noun. The package is attractive because its categorization and its Head relations fit together. Separate D can also license ordinary headedness, and determinative nouns can participate in fusion; neither possibility removes the reason to prefer the combination argued for here.

The package retains lexical and constructional restrictions. Nominal degree modifiers already occur in attributive AdjPs, so their distribution doesn't establish an added cost of NP projection. Coordinate rank alongside pronoun is a further preference: the comparisons here don't establish the additional shared rules that nesting would express.

\appendix

\section{Earlier accounts and logical alternatives}\label{sec:historical}

Keeping a word in one category across dependent and independent uses leaves its taxonomic position open. As \textcite[232--235]{lyons1968} observes, distribution can be compared at different levels: two expressions may belong together at one level and differ at a more specific level. Table~\ref{tab:rivals} distinguishes the selected authors' proposals, their analytical levels, and further logical alternatives.

\begin{table}[H]
\centering\small
\caption{Selected accounts and logical alternatives. The earlier proposals differ in scope and representational level; they don't all specify a surface taxonomy. The prose gives sources and inventory qualifications.}\label{tab:rivals}
\setlength{\tabcolsep}{4pt}
\begin{tabular}{>{\raggedright\arraybackslash}p{2.3cm}>{\raggedright\arraybackslash}p{6.2cm}>{\raggedright\arraybackslash}p{3.05cm}}
\toprule
Account & Proposed relationship & Level or status \\
\midrule
Palmer & Determinative adjectives grouped with pronouns. & Lexical grouping; inclusive Noun unspecified \\
Postal & Personal pronouns analysed as articles, with deeper noun features. & Intermediate and underlying representations \\
Sommerstein & Definite article and personal pronouns given NP structure. & Underlying representation \\
Lyons & Articles, demonstratives, and personal pronouns linked by definiteness and deixis. & Semantic relationship \\
\textit{CGEL} & Pronoun within Noun; determinative separate. & Lexical taxonomy \\
Hudson & Determiners within pronoun, which is within noun. & Lexical taxonomy \\
Spinillo & Determinatives redistributed; \mention{the}, \mention{a}, and \mention{every} retained as articles. & Lexical recategorization \\
D-noun analysis & Common noun, proper noun, pronoun, and determinative coordinate within Noun. & Proposed taxonomy \\
Logical alternative & Noun, pronoun, and determinative separate. & For comparison \\
Logical alternative & Determinative within Noun; pronoun separate. & For comparison \\
Logical alternative & Pronoun within determinative, which is within Noun. & For comparison \\
Logical alternative & Pronoun within proper noun; determinative within common noun. & For comparison \\
\bottomrule
\end{tabular}
\end{table}

\textcite[24]{palmer1924} proposed placing \enquote{determinative adjectives} with pronouns. He contrasted them with qualifying adjectives, which permit predicative use, comparison, and adverbial modification. Most determinatives, he observed, can \enquote{be used indifferently as pronouns or as modifiers of nouns}. \textcite[279]{lyons1968} adds a semantic connection: articles, demonstratives, and personal pronouns share definiteness and deictic contrasts.

\textcite{postal1966} develops the article--pronoun connection through English reflexives and combinations such as \mention{we men}. His category assignments depend on representational level: article segments occur at intermediate stages, their deepest counterparts are features of nouns, and surface pronouns can receive derivative Noun status. This differs from a proposal about the lexical taxonomy of the surface words.

\textcite[197--203]{sommerstein1972} argues in the opposite direction, giving the definite article and personal pronouns underlying NP structure. His English comparison includes the count restriction on anaphoric \mention{one}, whereas \mention{it} can refer to a quantity of a substance.\footnote{\textcite[797--798]{payne2013anaphoric} analyse anaphoric \mention{one} as a common count noun, distinct from determinative \mention{one} and personal pronoun \mention{one}.}

These predecessors challenge a fundamental article--pronoun separation but leave the full determinative inventory uncategorized. Its contrasts in grade, modification, complementation, and restricted independent use extend beyond their proposals. Hudson's explicit nesting and Spinillo's redistribution are assessed in §§\ref{sec:inheritance} and~\ref{sec:articles}, respectively.

\section*{Data and analysis materials}

The accompanying supplements are \href{run:matrix-audit.pdf}{\textit{Replication audit of the English determinative--pronoun feature matrix}} and \href{run:corpus-documentation.pdf}{\textit{CGELBank concordance and extraction notes}}. The \texttt{analysis/} directory preserves their input files, provenance records, scripts, numerical outputs, and sentence concordance. Its README identifies the files and reproduction procedures.

\vspace{0.5\baselineskip}
\noindent\textsc{Acknowledgements.} For the September 2026 revision, GPT-6 (Astra), Claude Opus~5, Claude Haiku~4.5, and GLM-5.3-Flash assisted drafting, source retrieval, script development, or critical review; GPT-5.6 (Sol) checked numerical outputs. The accompanying materials record the tools and their uses. Responsibility for the analysis is mine.

\clearpage
\printbibliography
\end{document}
